\clearpage
\phantomsection
\addcontentsline{toc}{chapter}{ABSTRACT}
\begin{center}
ABSTRACT

\vspace{0.25in}
ON THE NUMERICAL ANALYSIS OF MULTI-ROBOT SEARCH\\
AND RESCUE ALGORITHMS IN UNKNOWN\\
ENVIRONMENTS

\vspace{0.25in}
By

\vspace{0.1in}
Fozhan Babaeiyan Ghamsari

\vspace{0.1in}
Spring 2026
\end{center}

\begin{doublespacing}
\noindent
Every minute after a structural collapse, the probability of finding
survivors alive decreases. When autonomous robot teams are deployed to
search collapsed buildings, earthquake rubble, or flood sites, they
face two simultaneous constraints: a finite battery budget that forces
frequent returns to a charging base, and limited inter-robot
communication that makes coordination difficult.
This thesis studies the multi-robot search and rescue problem in which
$R$ robots must locate a target hidden among $n$ candidate sites, each
carrying a prior probability $p_i$ of containing the target, under a
finite per-robot energy budget $E$. Four algorithmic models spanning
the energy--communication design space are introduced and compared:
M1, an uncoordinated random baseline with \emph{no inter-robot
communication whatsoever}; M2, which introduces
\emph{Sequential Single-Item (SSI) auctions} with unconstrained energy;
M3, which adds the energy constraint with a standard $p_i/d$ bid
function; and M4/M4*, which extends each sortie via a \emph{greedy
chain extension} using the novel $p_i/d^2$ bid, alongside a
\emph{Hungarian assignment} baseline for comparison.
Performance is measured by the \emph{Finder Competitive Ratio} (FCR),
the ratio of the finding robot's travel distance to the omniscient
optimum, for which closed-form upper bounds are derived for each model.
All bounds are validated through 5,000 Monte Carlo trials in a
$10 \times 10$ square arena with $n = 30$ candidate sites and
$R = 3$ robots, with parameter sweeps over the energy budget
$E \in [8, 30]$ and fleet size $R \in [1, 6]$.
Four results emerge from this study.
Auction-based coordination yields a 6.6-fold improvement in FCR over
uncoordinated random search.
SSI auctions with probability-aware bids outperform Hungarian
assignment by $8.1\%$, and adding probability weighting to Hungarian
assignment alone \emph{worsens} performance due to a clustering
pathology, establishing that the SSI iterative bid-update mechanism
is essential.
Multi-node greedy chains recover $90\%$ of the performance gap between
energy-constrained and unconstrained operation.
The novel $p_i/d^2$ bid function outperforms the standard $p_i/d$ bid
by $1.5\%$ overall, with the advantage growing to approximately $12\%$
at tight energy budgets, justified by a cost-foreclosure argument
showing that a distant anchor site not only costs more energy to reach
but also shrinks the budget available for subsequent chain extensions.
Therefore, this thesis provides a reproducible foundational benchmark,
with closed-form guarantees and simple one-line deployment rules,
applicable to any scenario in which a team of resource-constrained
agents must locate a target under uncertainty---from urban
search-and-rescue to autonomous planetary exploration.
\end{doublespacing}
\clearpage
